
qui est un
foncteur sur la catégorie @@M0@@ (grâce à (IV 2.12)). En vertu de (IV 3.2.3), ce dernier commute aux limites inductives filtrantes. D'ailleurs, si @@M1@@ est une extension finie de @@M2@@, alors la projection @@M3@@ est un morphisme fini, plat, de présentation finie, donc il définit un homomorphisme image directe @@M4@@. On vérifie immédiatement que la condition envisagée dans 7.18.3 sur les homomorphismes norme est bien vérifiée. Enfin, si @@M5@@ est l'algèbre affine d'un ouvert @@M6@@ de l'espace affine @@M7@@, alors @@M8@@ s'identifie à l'application @@M9@@, qui est surjective pour @@M10@@ noethérien, en vertu de IX 1.6. Utilisant alors les résultats prouvés plus haut, on trouve :
\begin{enumerate}[label=(\roman*)]
\item Si @@M11@@ est une extension pure de @@M12@@, alors @@M13@@ est injectif, et même bijectif si @@M14@@ est noethérien.
\item Si @@M15@@ est séparablement clos et @@M16@@ une extension séparable, alors @@M17@@ est injectif.
\item Pour toute extension @@M18@@ de @@M19@@, le noyau de @@M20@@ est un groupe de torsion.
\end{enumerate}
Ces résultats restent valables si on suppose seulement @@M21@@ noethérien, à condition de se borner à des extensions @@M22@@ de @@M23@@ qui sont de type fini (de sorte que @@M24@@ est encore noethérien). De plus, on obtient essentiellement les mêmes résultats pour @@M25@@ au lieu du foncteur @@M26@@, à la seule différence près que dans (i), on ne peut affirmer en général qu'on a bijectivité, faute de disposer du théorème d'homotopie pour @@M27@@ au lieu de @@M28@@; pour pouvoir conclure à la bijectivité, il faudrait supposer @@M29@@ noethérien régulier, ce qui implique que les groupes @@M30@@ et @@M31@@ sont les mêmes pour @@M32@@, resp. pour @@M33@@.

